\section{Related work}
\label{sec:related}
\paragraph{Prediction-powered inference.} PPI~\citep{angelopoulos2023ppi} debiases the mean prediction on unlabeled data with a rectifier estimated on the labeled data, and gives CLT intervals for means, quantiles and M-estimation targets. PPI++~\citep{angelopoulos2023ppipp} adds a data-driven tuning weight $\lambda$ on the predictions, so that the interval is never asymptotically wider than the labeled-only one. Cross-prediction-powered inference~\citep{zrnic2023cross} treats the case where the predictor is trained on the same labels and uses cross-fitting. Active statistical inference~\citep{zrnic2024active} chooses which points to label from the model's uncertainty; stratified PPI~\citep{fisch2024stratified} uses strata of the prediction to allocate labels and to apply per-stratum rectifiers. Active selection needs control of the labeling mechanism and so is excluded by our fixed labeled set; the stratified idea (using $f$ non-linearly) is the closest in spirit to the calibration we study. The reference software is \texttt{ppi-python}~\citep{ppipython}.

\paragraph{Semiparametric efficiency and debiased learning.} Estimating a mean when the outcome is observed only for a random subset is the missing-data problem for which augmented inverse-probability-weighted estimators and their efficiency theory are classical~\citep{robins1994aipw,tsiatis2006}. PPI with predictor $f$ is the augmented estimator with the (generally misspecified) outcome regression $f$. Double/debiased machine learning~\citep{chernozhukov2018dml} supplies the cross-fitting device that lets a nuisance function be learned from the same sample without bias. Our estimators are this construction with a one-dimensional nuisance $g(f)$; what we add is the fixed-input question and its quantification on the standard PPI datasets, not a new estimator class.
